\subsection{GRADED SUBMODULES}

PROPOSITION 2. Let A be a graded ring of type \(\Delta\) , M a graded A-module of type \(\Delta\) , \((\mathbf{M}_{\lambda})\) its graduation and N a sub-A-module of M. The following properties are equivalent:

\begin{enumerate}
\def\labelenumi{(\alph{enumi})}
\item
  \(\mathbf{N}\) is the sum of the family \((\mathbf{N} \cap \mathbf{M}_{\lambda})_{\lambda \in \Delta}\) .
\item
  The homogeneous components of every element of N belong to N.
\item
  N is generated by homogeneous elements.
\end{enumerate}

Every element of N can be written uniquely as a sum of elements of the \(M_{\lambda}\) and hence it is immediate that (a) and (b) are equivalent and that (a) implies (c). We show that (c) implies (b). Then let \((x_{t})_{t\in I}\) be a family of homogeneous generators \(\neq0\) of N and let \(\delta(t)\) be the degree of \(x_{t}\) . Every element of N can be written as \(\sum_{i}a_{i}x_{i}\) with \(a_{i}\in A\) ; if \(a_{i,\lambda}\) is the component of \(a_{i}\) of degree \(\lambda\) , the conclusion follows from the relation

\[
\sum_ {t \in I} \left(\sum_ {\mu \in \Delta} a _ {t, \mu} x _ {t}\right) = \sum_ {\lambda \in \Delta} \left(\sum_ {\mu + \delta (t) = \lambda} a _ {t, \mu} x _ {t}\right).
\]

Remark (1) In the above notation, the relation \(\sum_{t\in I}a_t x_t = 0\) is therefore equivalent to the system of relations \(\sum_{\mu +\delta (t) = \lambda}a_{t,\mu}x_{t} = 0\) . When \(\Delta\) is a group, these relations can be written \(\sum_{t\in I}a_{t,\lambda -\delta (t)}x_{t} = 0\) .

When a submodule N of M has the equivalent properties stated in Proposition 2, clearly the \(N \cap M_{\lambda}\) form a graduation compatible with the A-module structure of N, called the graduation induced by that on M; N with this graduation is called a graded submodule of M.

COROLLARY 1. If N is a graded submodule of M and \((x_{t})\) is a generating system of N, the homogeneous components of the \(x_{t}\) form a generating system of N.

COROLLARY 2. If N is a finitely generated submodule of M, N admits a finite generating system consisting of homogeneous elements.

It suffices to apply Corollary 1 noting that an element of M has only a finite number of homogeneous components \(\neq0\) .

A graded submodule of \(A_{s}\) (resp. \(A_{a}\) ) is called a graded left (resp. right) ideal of the graded ring A. For every subring B of \(\mathrm{A}(\mathrm{B} \cap \mathrm{A}_{\lambda})(\mathrm{B} \cap \mathrm{A}_{\mu}) \subset \mathrm{B} \cap \mathrm{A}_{\lambda + \mu}\) ; if B is a graded sub-Z-module of A, the graduation induced on B by that on A is therefore compatible with the ring structure on B; B is then called a graded subring of A.

Clearly if N (resp. B) is a graded sub-A-module of M (resp. a graded subring of A), the canonical injection \(N \rightarrow M\) (resp. \(B \rightarrow A\) ) is a graded module homomorphism of degree 0 (resp. a graded ring homomorphism).

If N is a graded submodule of a graded A-module M and \((\mathbf{M}_{\lambda})_{\lambda\in\Delta}\) the graduation of M, the submodules \((\mathbf{M}_{\lambda}+\mathbf{N})/\mathbf{N}\) of M/N form a graduation compatible with the structure of this quotient module. For, if \(N_{\lambda}=M_{\lambda}\cap N\) , \((\mathbf{M}_{\lambda}+\mathbf{N})/\mathbf{N}\) is identified with \(M_{\lambda}/N_{\lambda}\) and it follows from Proposition 2 and §1, no. 6, formula (26) that M/N is their direct sum. Moreover,

\[
\mathrm{A} _ {\lambda} (\mathrm{M} _ {\mu} + \mathrm{N}) \subset \mathrm{A} _ {\lambda} \mathrm{M} _ {\mu} + \mathrm{N} \subset \mathrm{M} _ {\lambda + \mu} + \mathrm{N}
\]

and hence \(\mathrm{A}_{\lambda}((\mathrm{M}_{\mu}+\mathrm{N})/\mathrm{N})\subset(\mathrm{M}_{\lambda+\mu}+\mathrm{N})/\mathrm{N}\) , which establishes our assertion. The graduation \(((\mathbf{M}_{\lambda} + \mathbf{N})/\mathbf{N})_{\lambda \in \Delta}\) is called the quotient graduation of that on M by N and the quotient module M/N with this graduation is called the graded quotient module of M by the graded submodule N; the canonical homomorphism \(M \to M/N\) is a graded homomorphism of degree 0 for this graduation.

If b is a graded two-sided ideal of A, the quotient graduation on A/b is compatible with the ring structure on A/b; the ring A/b with this graduation is called the quotient graded ring of A by b; the canonical homomorphism \(A \rightarrow A/b\) is a homomorphism of graded rings for this graduation.

PROPOSITION 3. Let A be a graded ring of type \(\Delta\) , M, N two graded A-modules of type \(\Delta\) and \(u\) : M \(\rightarrow\) N a graded A-homomorphism of degree \(\delta\) . Then:

\begin{enumerate}
\def\labelenumi{(\roman{enumi})}
\item
  \(\operatorname{Im}(u)\) is a graded submodule of \(\mathbf{N}\) .
\item
  If \(\delta\) is a regular element of \(\Delta\) , \(\operatorname{Ker}(u)\) is a graded submodule of \(\mathbf{M}\) .
\item
  If \(\delta = 0\) , the bijection \(\mathbf{M} / \mathrm{Ker}(u) \to \operatorname{Im}(u)\) canonically associated with \(u\) is an isomorphism of graded modules.
\end{enumerate}

Assertion (i) follows immediately from the definitions and Proposition 2(c). If \(x\) is an element of \(\mathbf{M}\) such that \(u(x) = 0\) and \(x = \sum_{\lambda} x_{\lambda}\) is its decomposition into homogeneous components (where \(x_{\lambda}\) is of degree \(\lambda\) ), then

\[
\sum_ {\lambda} u (x _ {\lambda}) = u (x) = 0
\]

and \(u(x_{\lambda})\) is of degree \(\lambda + \delta\) ; if \(\delta\) is regular the relation \(\lambda + \delta = \mu + \delta\) implies \(\lambda = \mu\) , hence the \(u(x_{\lambda})\) are the homogeneous components of \(u(x)\) and necessarily \(u(x_{\lambda}) = 0\) for all \(\lambda \in \Delta\) , which proves (ii). The bijection \(v: \mathbf{M} / \operatorname{Ker}(u) \to \operatorname{Im}(u)\) canonically associated with \(u\) is then a graded homomorphism of degree \(\delta\) , as follows from the definition of the quotient graduation; whence (iii) when \(\delta = 0\) .

COROLLARY. Let A, B be two graded rings of type \(\Delta\) and u: A \(\rightarrow\) B a graded homomorphism of graded rings. Then \(\operatorname{Im}(u)\) is a graded submodule of B, \(\operatorname{Ker}(u)\) a graded two-sided ideal of A and the bijection A/Ker(u) \(\rightarrow\) \(\operatorname{Im}(u)\) canonically associated with u is an isomorphism of graded rings.

It suffices to apply Proposition 3 to u considered as a homomorphism of degree 0 of graded Z-modules.

PROPOSITION 4. Let A be a graded ring of type \(\Delta\) and M a graded A-module of type \(\Delta\) .

\begin{enumerate}
\def\labelenumi{(\roman{enumi})}
\item
  Every sum and every intersection of graded submodules of \(\mathbf{M}\) is a graded submodule.
\item
  If \(x\) is a homogeneous element of \(\mathbf{M}\) of degree \(\mu\) which is cancellable in \(\Delta\) , the annihilator of \(x\) is a graded left ideal of \(\mathbf{A}\) .
\item
  If all the elements of \(\Delta\) are cancellable, the annihilator of a graded submodule of M is a graded two-sided ideal of A.
\end{enumerate}

If \((\mathbf{N}_{t})\) is a family of graded submodules of M, property (c) of Proposition 2 shows that the sum of the \(N_{t}\) is generated by homogeneous elements and property (b) of Proposition 2 proves that the homogeneous components of every element of \(\bigcap_{t} N_{t}\) belongs to \(\bigcap_{t} N_{t}\) ; whence (i).

To prove (ii), it suffices to note that \(\operatorname{Ann}(x)\) is the kernel of the homomorphism \(a \mapsto ax\) of the A-module \(\mathbf{A}_s\) into \(\mathbf{M}\) and that this homomorphism is graded of degree \(\mu\) ; the conclusion follows from Proposition 3(ii). Finally (iii) is a consequence of (i) and (ii) for the annihilator of a graded submodule \(\mathbf{N}\) of \(\mathbf{M}\) is the intersection of the annihilators of the homogeneous elements of \(\mathbf{N}\) , by virtue of Proposition 2.

Remark 2. Let M be a graded A-module and E a submodule of M; it follows from Proposition 4(i) that there exists a largest graded submodule N' of M contained in E and a smallest graded submodule N'' of M containing E; N' is the set of x \ensuremath{\in} E all of whose homogeneous components belong to E and N'' is the submodule of M generated by the homogeneous components of a generating system of E.

PROPOSITION 5. Let A be a graded ring of type \(\Delta\) . If every element of \(\Delta\) is cancellable, then, for every homogeneous element \(a \in \mathbf{A}\) , the centralizer of \(a\) in A (I, § 1, no. 5) is a graded subring of A.

Suppose that \(a\) is of degree \(\delta\) ; let \(b = \sum_{\lambda} b_{\lambda}\) be an element permutable with \(a, b_{\lambda}\) being the homogeneous component of \(b\) of degree \(\lambda\) for all \(\lambda \in \Delta\) . Then by hypothesis \(\sum_{\lambda} (ab_{\lambda} - b_{\lambda}a) = 0\) and \(ab_{\lambda} - b_{\lambda}a\) is homogeneous of degree \(\lambda + \delta\) ; as \(\delta\) is cancellable, it follows that \(ab_{\lambda} = b_{\lambda}a\) for all \(\lambda\) , which proves our assertion.

COROLLARY. If every element of \(\Delta\) is cancellable, the centralizer of the graded subring B of A (and in particular the centre of A) is a graded subring of A.

It is the intersection of the centralizers of the homogeneous elements of B.

Remark (3) A direct system \((\mathbf{A}_{\alpha}, \phi_{\beta\alpha})\) of graded rings of type \(\Delta\) (resp. a direct system \((\mathbf{M}_{\alpha}, f_{\beta\alpha})\) of graded \(A_{\alpha}\) -modules of type \(\Delta\) ) is a direct system of rings (resp. \(A_{\alpha}\) -modules) such that each \(A_{\alpha}\) (resp. \(M_{\alpha}\) ) is graded of type \(\Delta\) and each \(\phi_{\beta\alpha}\) (resp. \(f_{\beta\alpha}\) ) is a homomorphism of graded rings (resp. an \(A_{\alpha}\) -homomorphism of degree 0 of graded modules). If \((\mathbf{A}_{\alpha}^{\lambda})_{\lambda \in \Delta}\) (resp. \((\mathbf{M}_{\alpha}^{\lambda})_{\lambda \in \Delta}\) ) be the graduation of \(A_{\alpha}\) (resp. \(M_{\alpha}\) ) and we write

\[
\mathrm{A} = \lim _ {\rightarrow} \mathrm{A} _ {\alpha}, \quad \mathrm{A} ^ {\lambda} = \lim _ {\overrightarrow {\alpha}} \mathrm{A} _ {\alpha} ^ {\lambda} \quad (\text { resp.   } \mathrm{M} = \lim _ {\rightarrow} \mathrm{M} _ {\alpha}, \mathrm{M} ^ {\lambda} = \lim _ {\overrightarrow {\alpha}} \mathrm{M} _ {\alpha} ^ {\lambda}),
\]

it follows from § 6, no. 2, Proposition 5 that \((\mathbf{A}^{\wedge})\) (resp. \((\mathbf{M}^{\wedge}))\) is a graduation of A (resp. M) and it follows from I, § 10, nos. 3 and 4 that this graduation is compatible with the ring structure on A (resp. the A-module structure on M). The graded ring A (resp. graded A-module M) is called the direct limit of the direct system of graded rings \((\mathbf{A}_{\alpha}, \phi_{\beta \alpha})\) (resp. graded modules \((\mathbf{M}_{\alpha}, f_{\beta \alpha})\) ). If \(\phi_{\alpha} \colon \mathbf{A}_{\alpha} \to \mathbf{A}\) (resp. \(f_{\alpha} \colon \mathbf{M}_{\alpha} \to \mathbf{M}\) ) is the canonical mapping, \(\phi_{\alpha}\) (resp. \(f_{\alpha}\) ) is a homomorphism of graded rings (resp. a homomorphism of degree 0 of graded \(\mathbf{A}_{\alpha}\) -modules).

